% Figure: a two-point climbing anchor. Two placements share a load through a
% sling to a shared master point; the included angle sets the leg tension.
% (a) the anchor geometry; (b) leg tension vs included angle.
\begin{figure}[htbp]
\centering
\begin{tikzpicture}[
  force/.style={draw=engblue, -{Stealth}, very thick},
  rxn/.style={draw=engred, -{Stealth}, very thick},
  sling/.style={draw=engamber, line width=1.4pt},
  scale=1.0,
]
% ===== (a) two-point anchor =====
\begin{scope}
  % rock face at top
  \draw[draw=engblack, very thick] (-0.4,3.2) -- (3.8,3.2);
  \foreach \x in {-0.3,0.1,...,3.7} {\draw[draw=enggray, thin] (\x,3.2) -- (\x-0.2,3.46);}
  % two placements
  \coordinate (P1) at (0.4,3.2);
  \coordinate (P2) at (3.0,3.2);
  \fill[engblack] (P1) circle (0.07);
  \fill[engblack] (P2) circle (0.07);
  \node[font=\scriptsize, anchor=south] at (P1) {placement 1};
  \node[font=\scriptsize, anchor=south] at (P2) {placement 2};
  % master point below
  \coordinate (M) at (1.7,1.2);
  \draw[sling] (P1) -- (M);
  \draw[sling] (P2) -- (M);
  \draw[draw=engblack, fill=enggray!20, thick] (M) circle (0.11);
  % load down at master point
  \draw[force] (M) -- ++(0,-1.4) node[anchor=north, font=\small, engblue] {$W$};
  % leg tensions
  \draw[rxn] (M) -- ($(M)!0.42!(P1)$) node[anchor=south east, font=\scriptsize, engred] {$T_1$};
  \draw[rxn] (M) -- ($(M)!0.42!(P2)$) node[anchor=south west, font=\scriptsize, engred] {$T_2$};
  % included angle
  \draw[draw=enggray, thin] ($(M)!0.7!(P1)$) arc (128:52:0.85);
  \node[font=\scriptsize, enggray] at (1.7,2.05) {$2\theta$};
  \node[font=\small, anchor=north] at (1.7,-0.4) {(a) two-point anchor};
\end{scope}
% ===== (b) leg tension vs included angle =====
\begin{scope}[xshift=6.4cm, yshift=0.1cm]
  \begin{axis}[
    width=6.4cm, height=5.2cm,
    xlabel={included angle $2\theta$ (deg)},
    ylabel={leg tension $/\,W$},
    xmin=0, xmax=170, ymin=0.4, ymax=3.2,
    xtick={0,60,120,180},
    ytick={0.5,1,1.5,2,2.5,3},
    grid=both, grid style={enggray!25},
    label style={font=\scriptsize},
    tick label style={font=\scriptsize},
    every axis plot/.append style={line width=1.2pt},
  ]
    % T/W = 1/(2 cos theta), theta = x/2 degrees
    \addplot[draw=engred, domain=0:165, samples=120] {1/(2*cos(x/2))};
    \node[font=\scriptsize, anchor=west, engred] at (axis cs:10,0.65) {$T/W = \tfrac{1}{2\cos\theta}$};
  \end{axis}
  \node[font=\small, anchor=north] at (3.0,-0.4) {(b) the angle penalty};
\end{scope}
\end{tikzpicture}
\caption[A two-point climbing anchor and its angle penalty]{(a) Two rock
placements share a load through a sling to a single master point, the fitting
a climber clips into. The two legs meet at the master point with an included
angle $2\theta$. Each leg carries a tension $T = W/(2\cos\theta)$, the same
hanging-mass result read as a rig the climber builds. (b) The leg tension per
unit load climbs slowly to $60^{\circ}$ of included angle, then sharply: at a
$120^{\circ}$ included angle each leg already carries the full load, and past
that the tension runs away. The working rule that an anchor's included angle
should stay under about $60^{\circ}$ is this curve read for its knee, and a
wide sling on a narrow ledge is the field mistake that overloads both
placements at once.}
\label{fig:vol03ch01:climbing-anchor}
\end{figure}
